\chapter[Convex Duality, KKT Conditions \& SVMs]{Convex Duality, KKT Conditions \& Support Vector Machines}
\label{chap:week10}

% ----------------------------------------------------------------------
% PART 1: EXECUTIVE ROADMAP & ARCHITECTURAL FOUNDATIONS
% ----------------------------------------------------------------------
\section{Executive Roadmap \& Learning Objectives}

\subsection{Opening Hook: The Infinite Void and the Mountain Pillars}
Imagine an artist tasked with sculpting an infinite-dimensional mountain range ($d = \infty$). Every microscopic fold, crevice, and ridge must be positioned with mathematical perfection. If the artist is forced to define the mountain by specifying every infinitesimal coordinate, the task is utterly impossible: no human lifetime or computer memory could ever store an infinite sequence of weights.

Now imagine that a geometer enters the studio and reveals a profound structural secret: the entire soaring mountain range is supported from beneath by just three immovable pillars of bedrock. If you know the exact location and upward thrust of those three pillars, the entire mountain's shape is mathematically locked into place. You do not need to sculpt the infinite void of space; you only need to determine the mechanical forces acting at the three contact points.

This is the miracle of \textbf{Lagrange Duality} and \textbf{Support Vector Machines (SVMs)}. 

When machine learning algorithms project data into infinite-dimensional Reproducing Kernel Hilbert Spaces (RKHS) via nonlinear kernels (such as the Gaussian RBF kernel), the primal weight vector $\mathbf{w}$ contains infinitely many parameters. Attempting to optimize $\mathbf{w}$ directly in primal space is computationally intractable. 

Lagrange duality inverts the entire dimensionality contract: instead of optimizing over infinitely many feature dimensions $d = \infty$, the dual problem optimizes over exactly $N$ scalar multipliers $\boldsymbol{\lambda} \in \mathbb{R}^N$---one scalar per training sample. Even more remarkably, the Karush-Kuhn-Tucker (KKT) complementary slackness conditions guarantee that for every data point residing comfortably inside the class territory, its dual multiplier is identically zero ($\lambda_i^* = 0$). 

The infinite-dimensional decision boundary collapses into a sparse linear combination of just a tiny handful of critical boundary samples: the \textbf{support vectors}. The boundary does not care about the millions of points in the interior; it is mechanically pinned by the bedrock samples on the margin.

\subsection{Architectural Bridge from Previous Chapters}
In Chapter~\ref{chap:week09}, we established the geometric foundation of constrained continuous optimization, analyzing Lagrange multipliers for equality-constrained manifolds and proving the non-expansiveness of Euclidean projection onto closed convex sets. 

However, practical machine learning classification rarely involves simple equality constraints. Class boundaries are defined by \textbf{inequalities}: we require samples to lie on or beyond a decision margin:
\begin{equation}
y_i(\mathbf{w}^T \mathbf{x}_i + b) \ge 1 \iff 1 - y_i(\mathbf{w}^T \mathbf{x}_i + b) \le 0
\end{equation}
Chapter 10 generalizes the Lagrangian apparatus to arbitrary systems of convex inequality and affine equality constraints. We formulate the \textbf{Lagrange Dual Problem}, prove weak and strong duality theorems, develop the complete four-part \textbf{Karush-Kuhn-Tucker (KKT) optimality conditions}, and derive the Hard-Margin and Soft-Margin Support Vector Machine. This chapter transforms constrained optimization from a theoretical framework into one of the most powerful and robust classification algorithms in machine learning.

\subsection{Prerequisites \& Scaffolding}
Mastery of this chapter synthesizes convex functions, hyperplanes, and multivariable calculus as outlined in Table~\ref{tab:w10_prereqs}.

\begin{table}[htbp]
\centering
\small
\begin{tabularx}{\textwidth}{lXX}
\toprule
\textbf{Foundational Area} & \textbf{Core Concepts Required} & \textbf{Role in Machine Learning Foundations} \\
\midrule
Convex Analysis & Convex sets, epigraphs, supporting hyperplanes, Slater's condition & Guaranteeing zero duality gap and establishing necessary and sufficient KKT conditions \\
\addlinespace
Vector Calculus & Multivariable gradients, partial derivatives, stationarity & Formulating the Lagrangian stationarity conditions $\nabla_{\mathbf{w}}\mathcal{L} = \mathbf{0}$ \\
\addlinespace
Linear Algebra & Inner products, Gram kernel matrices $K = X X^T$, positive definiteness & Deriving the dual Quadratic Program (QP) and unlocking the kernel trick \\
\bottomrule
\end{tabularx}
\caption{Prerequisite Scaffolding for Chapter 10.}
\label{tab:w10_prereqs}
\end{table}

\subsection{Core Competencies}
By the conclusion of this chapter, students will be able to:
\begin{itemize}
    \item Formulate the Lagrange dual function $g(\boldsymbol{\lambda}, \boldsymbol{\nu})$ and prove that it is universally concave regardless of primal convexity.
    \item State and verify Weak Duality, Strong Duality, and Slater's constraint qualification.
    \item Formulate and solve the complete four-part Karush-Kuhn-Tucker (KKT) optimality conditions.
    \item Rigorously derive the Hard-Margin and Soft-Margin Support Vector Machine from geometric margin maximization.
    \item Apply complementary slackness to prove why non-support vectors exert zero influence on the decision boundary.
    \item Derive the box constraints $0 \le \lambda_i \le C$ and classify support vectors into free margin points and bounded margin violators.
    \item Formulate the Dual Quadratic Program and implement an SVM solver using scientific computing libraries.
\end{itemize}

\subsection{Notation \& Mathematical Conventions}
Throughout this chapter, we adhere to the standard mathematical conventions:
\begin{itemize}
    \item Primal objective functional: $f_0(\mathbf{w}): \mathbb{R}^d \to \mathbb{R}$.
    \item Inequality constraint functions: $f_i(\mathbf{w}) \le 0$ ($i = 1, \dots, m$).
    \item Equality constraint functions: $h_j(\mathbf{w}) = 0$ ($j = 1, \dots, p$).
    \item Lagrangian functional: $\mathcal{L}(\mathbf{w}, \boldsymbol{\lambda}, \boldsymbol{\nu}) \triangleq f_0(\mathbf{w}) + \sum_{i=1}^m \lambda_i f_i(\mathbf{w}) + \sum_{j=1}^p \nu_j h_j(\mathbf{w})$.
    \item Lagrange dual function: $g(\boldsymbol{\lambda}, \boldsymbol{\nu}) \triangleq \inf_{\mathbf{w} \in \mathbb{R}^d} \mathcal{L}(\mathbf{w}, \boldsymbol{\lambda}, \boldsymbol{\nu})$.
    \item Primal optimal value: $p^* \triangleq \inf_{\mathbf{w} \in \text{feasible}} f_0(\mathbf{w})$.
    \item Dual optimal value: $d^* \triangleq \sup_{\boldsymbol{\lambda} \ge \mathbf{0}, \boldsymbol{\nu}} g(\boldsymbol{\lambda}, \boldsymbol{\nu})$.
    \item Binary class labels: $y_i \in \{-1, +1\}$.
    \item Slack variables: $\xi_i \ge 0$, measuring margin violations in Soft-Margin SVM.
    \item Regularization penalty parameter: $C > 0$, balancing margin width against margin slack penalties.
\end{itemize}

% ----------------------------------------------------------------------
% PART 2: MATHEMATICAL & COMPUTATIONAL FOUNDATIONS
% ----------------------------------------------------------------------
\section{Mathematical \& Computational Foundations}

\subsection{The Primal and Dual Problems}
Consider the canonical constrained continuous optimization problem in standard form:
\begin{equation}
\min_{\mathbf{w} \in \mathbb{R}^d} f_0(\mathbf{w}) \quad \text{subject to } \begin{cases} f_i(\mathbf{w}) \le 0, & i = 1, \dots, m \\ h_j(\mathbf{w}) = 0, & j = 1, \dots, p \end{cases}
\label{eq:standard_primal}
\end{equation}
with optimal value $p^* = \inf \{f_0(\mathbf{w}) \mid f_i(\mathbf{w}) \le 0, \, h_j(\mathbf{w}) = 0\}$.

\begin{definitionbox}{The Lagrange Dual Function}
The \textbf{Lagrange Dual Function} $g: \mathbb{R}^m \times \mathbb{R}^p \to \mathbb{R} \cup \{-\infty\}$ is the infimum of the Lagrangian over the primal parameter vector $\mathbf{w} \in \mathbb{R}^d$:
\begin{equation}
g(\boldsymbol{\lambda}, \boldsymbol{\nu}) \triangleq \inf_{\mathbf{w} \in \mathbb{R}^d} \left( f_0(\mathbf{w}) + \sum_{i=1}^m \lambda_i f_i(\mathbf{w}) + \sum_{j=1}^p \nu_j h_j(\mathbf{w}) \right)
\label{eq:dual_func_def}
\end{equation}
where $\boldsymbol{\lambda} \in \mathbb{R}^m$ is the dual multiplier vector associated with the inequality constraints, and $\boldsymbol{\nu} \in \mathbb{R}^p$ is the dual multiplier vector associated with the equality constraints.
\end{definitionbox}

\begin{theorembox}{Universal Concavity of the Dual Function}
The Lagrange dual function $g(\boldsymbol{\lambda}, \boldsymbol{\nu})$ is always concave in $(\boldsymbol{\lambda}, \boldsymbol{\nu})$, even if the primal objective $f_0$ and the constraint functions $f_i, h_j$ are wildly non-convex.
\end{theorembox}

\begin{proof}
For any fixed parameter vector $\mathbf{w} \in \mathbb{R}^d$, define the affine function $\psi_{\mathbf{w}}: \mathbb{R}^m \times \mathbb{R}^p \to \mathbb{R}$:
\begin{equation}
\psi_{\mathbf{w}}(\boldsymbol{\lambda}, \boldsymbol{\nu}) \triangleq f_0(\mathbf{w}) + \sum_{i=1}^m \lambda_i f_i(\mathbf{w}) + \sum_{j=1}^p \nu_j h_j(\mathbf{w})
\end{equation}
Notice that for a fixed $\mathbf{w}$, $\psi_{\mathbf{w}}$ is strictly linear (affine) with respect to the dual multipliers $(\boldsymbol{\lambda}, \boldsymbol{\nu})$.
By definition, the dual function is the pointwise infimum of this family of affine functions:
\begin{equation}
g(\boldsymbol{\lambda}, \boldsymbol{\nu}) = \inf_{\mathbf{w} \in \mathbb{R}^d} \psi_{\mathbf{w}}(\boldsymbol{\lambda}, \boldsymbol{\nu})
\end{equation}
Let $\theta \in [0, 1]$ and consider two arbitrary dual points $(\boldsymbol{\lambda}_1, \boldsymbol{\nu}_1)$ and $(\boldsymbol{\lambda}_2, \boldsymbol{\nu}_2)$:
\begin{align}
g(\theta \boldsymbol{\lambda}_1 + (1-\theta)\boldsymbol{\lambda}_2, \theta \boldsymbol{\nu}_1 + (1-\theta)\boldsymbol{\nu}_2)
&= \inf_{\mathbf{w}} \psi_{\mathbf{w}}(\theta \boldsymbol{\lambda}_1 + (1-\theta)\boldsymbol{\lambda}_2, \theta \boldsymbol{\nu}_1 + (1-\theta)\boldsymbol{\nu}_2) \nonumber \\
&= \inf_{\mathbf{w}} \left[ \theta \psi_{\mathbf{w}}(\boldsymbol{\lambda}_1, \boldsymbol{\nu}_1) + (1-\theta) \psi_{\mathbf{w}}(\boldsymbol{\lambda}_2, \boldsymbol{\nu}_2) \right] \nonumber \\
&\ge \theta \inf_{\mathbf{w}} \psi_{\mathbf{w}}(\boldsymbol{\lambda}_1, \boldsymbol{\nu}_1) + (1-\theta) \inf_{\mathbf{w}'} \psi_{\mathbf{w}'}(\boldsymbol{\lambda}_2, \boldsymbol{\nu}_2) \nonumber \\
&= \theta g(\boldsymbol{\lambda}_1, \boldsymbol{\nu}_1) + (1-\theta) g(\boldsymbol{\lambda}_2, \boldsymbol{\nu}_2)
\end{align}
Because the secant chord lies everywhere below the dual function graph, $g(\boldsymbol{\lambda}, \boldsymbol{\nu})$ is universally concave.
\end{proof}

\begin{theorembox}{Weak Duality Theorem}
For any dual feasible multipliers $\boldsymbol{\lambda} \ge \mathbf{0}$ and $\boldsymbol{\nu} \in \mathbb{R}^p$, the dual function value provides a universal lower bound on the primal optimal value:
\begin{equation}
g(\boldsymbol{\lambda}, \boldsymbol{\nu}) \le p^*
\label{eq:weak_duality}
\end{equation}
Consequently, the optimal dual value $d^* \triangleq \sup_{\boldsymbol{\lambda} \ge \mathbf{0}, \boldsymbol{\nu}} g(\boldsymbol{\lambda}, \boldsymbol{\nu})$ satisfies:
\begin{equation}
d^* \le p^*
\end{equation}
The non-negative difference $p^* - d^* \ge 0$ is the \textbf{Duality Gap}.
\end{theorembox}

\begin{proof}
Let $\tilde{\mathbf{w}}$ be any feasible point for the primal problem, meaning $f_i(\tilde{\mathbf{w}}) \le 0$ for all $i$ and $h_j(\tilde{\mathbf{w}}) = 0$ for all $j$.
Because $\lambda_i \ge 0$ and $f_i(\tilde{\mathbf{w}}) \le 0$, the product $\lambda_i f_i(\tilde{\mathbf{w}}) \le 0$. Thus:
\begin{equation}
\mathcal{L}(\tilde{\mathbf{w}}, \boldsymbol{\lambda}, \boldsymbol{\nu}) = f_0(\tilde{\mathbf{w}}) + \sum_{i=1}^m \underbrace{\lambda_i f_i(\tilde{\mathbf{w}})}_{\le 0} + \sum_{j=1}^p \underbrace{\nu_j h_j(\tilde{\mathbf{w}})}_{= 0} \le f_0(\tilde{\mathbf{w}})
\end{equation}
By definition of the infimum over all $\mathbf{w} \in \mathbb{R}^d$:
\begin{equation}
g(\boldsymbol{\lambda}, \boldsymbol{\nu}) = \inf_{\mathbf{w}} \mathcal{L}(\mathbf{w}, \boldsymbol{\lambda}, \boldsymbol{\nu}) \le \mathcal{L}(\tilde{\mathbf{w}}, \boldsymbol{\lambda}, \boldsymbol{\nu}) \le f_0(\tilde{\mathbf{w}})
\end{equation}
Because $g(\boldsymbol{\lambda}, \boldsymbol{\nu}) \le f_0(\tilde{\mathbf{w}})$ holds for \emph{every} primal feasible point $\tilde{\mathbf{w}}$, taking the infimum over all feasible $\tilde{\mathbf{w}}$ establishes $g(\boldsymbol{\lambda}, \boldsymbol{\nu}) \le p^*$.
\end{proof}

\begin{ponderbox}{Why Must the Dual Problem Always Be Concave?}
Why is the universal concavity of the dual function $g(\boldsymbol{\lambda}, \boldsymbol{\nu})$ so profound for continuous optimization?
\tcblower
\textbf{The Socratic Resolution:} Maximizing a concave function over a convex set (the non-negative orthant $\boldsymbol{\lambda} \ge \mathbf{0}$) is, by definition, a \textbf{convex optimization problem}! 

Even if the primal problem $f_0(\mathbf{w})$ is wildly non-convex, multimodal, discontinuous, or NP-hard, its Lagrange dual problem $\max_{\boldsymbol{\lambda} \ge \mathbf{0}, \boldsymbol{\nu}} g(\boldsymbol{\lambda}, \boldsymbol{\nu})$ is \emph{always} a well-behaved convex program without spurious local maxima. Solving the dual problem via subgradient ascent or interior-point methods produces a certifiable, non-trivial global lower bound $d^* \le p^*$ on the unknown global minimum.
\end{ponderbox}

\subsection{Strong Duality and the KKT Optimality Conditions}

\begin{definitionbox}{Slater's Constraint Qualification}
Let the primal problem~\eqref{eq:standard_primal} be convex: the objective $f_0$ and inequality constraints $f_1, \dots, f_m$ are convex, and equality constraints $h_j(\mathbf{w}) = \mathbf{a}_j^T \mathbf{w} - b_j$ are affine.
If there exists a \textbf{strictly feasible point} $\tilde{\mathbf{w}} \in \mathbb{R}^d$ such that:
\begin{equation}
f_i(\tilde{\mathbf{w}}) < 0 \quad (\forall i \in \{1, \dots, m\}) \quad \text{and} \quad A\tilde{\mathbf{w}} = \mathbf{b}
\label{eq:slaters_condition}
\end{equation}
then \textbf{Strong Duality} holds: the duality gap is zero ($d^* = p^*$), and the dual supremum is attained at some dual optimal $(\boldsymbol{\lambda}^*, \boldsymbol{\nu}^*)$.
\end{definitionbox}

\begin{theorembox}{Karush-Kuhn-Tucker (KKT) Optimality Conditions}
Let $f_0, f_1, \dots, f_m$ be continuously differentiable convex functions, and let $h_j$ be affine. If Slater's condition holds, then $\mathbf{w}^*$ is primal optimal and $(\boldsymbol{\lambda}^*, \boldsymbol{\nu}^*)$ is dual optimal if and only if they satisfy the complete four-part KKT conditions:
\begin{enumerate}
    \item \textbf{Primal Feasibility:}
    \begin{align}
    f_i(\mathbf{w}^*) &\le 0 \quad (i = 1, \dots, m) \\
    h_j(\mathbf{w}^*) &= 0 \quad (j = 1, \dots, p)
    \end{align}
    \item \textbf{Dual Feasibility:}
    \begin{equation}
    \lambda_i^* \ge 0 \quad (i = 1, \dots, m)
    \end{equation}
    \item \textbf{Complementary Slackness:}
    \begin{equation}
    \lambda_i^* f_i(\mathbf{w}^*) = 0 \quad (i = 1, \dots, m)
    \label{eq:comp_slackness}
    \end{equation}
    \item \textbf{Lagrangian Stationarity:}
    \begin{equation}
    \nabla_{\mathbf{w}} \mathcal{L}(\mathbf{w}^*, \boldsymbol{\lambda}^*, \boldsymbol{\nu}^*) = \nabla f_0(\mathbf{w}^*) + \sum_{i=1}^m \lambda_i^* \nabla f_i(\mathbf{w}^*) + \sum_{j=1}^p \nu_j^* \nabla h_j(\mathbf{w}^*) = \mathbf{0}
    \label{eq:kkt_stationarity}
    \end{equation}
\end{enumerate}
\end{theorembox}

\begin{ponderbox}{Geometric Anatomy of the Three Multiplier Regimes}
In a constrained optimization problem with inequality constraints $f_i(\mathbf{w}) \le 0$, how does the complementary slackness condition $\lambda_i^* f_i(\mathbf{w}^*) = 0$ act as a logical and geometric switch?
\tcblower
\textbf{The Socratic Resolution:} Because $\lambda_i^* \ge 0$ and $f_i(\mathbf{w}^*) \le 0$, their product vanishes if and only if at least one of the two factors is zero. This partitions all training points into two mutually exclusive geometric regimes:
\begin{itemize}
    \item \textbf{Strictly Inactive Regime ($f_i(\mathbf{w}^*) < 0 \implies \lambda_i^* = 0$):} The point lies in the interior of the feasible set, away from the boundary. Tightening or removing this constraint locally would not change the optimal solution. The constraint exerts zero contact force: $\lambda_i^* = 0$.
    \item \textbf{Active Boundary Regime ($\lambda_i^* > 0 \implies f_i(\mathbf{w}^*) = 0$):} The point presses directly against the constraint boundary. If the boundary were relaxed, the optimal solution would move into the newly opened territory. The multiplier $\lambda_i^* > 0$ represents the exact normal reaction force required to prevent the solution from drifting across the boundary.
\end{itemize}
\end{ponderbox}

% ----------------------------------------------------------------------
% PART 3: VISUAL & GEOMETRIC ANCHORING
% ----------------------------------------------------------------------
\section{Visual \& Geometric Anchoring}

In this section, we anchor the mathematical apparatus of duality and SVMs in three native TikZ visualizations.

\subsection{The Geometry of Support Vector Machines and Margins}
Figure~\ref{fig:svm_geometry} illustrates the canonical maximum margin hyperplane separating two linearly separable classes. The decision boundary is determined strictly by the circled support vectors that lie on the margin boundaries; all interior points have $\lambda_i^* = 0$.

\begin{figure}[htbp]
\centering
\begin{tikzpicture}[scale=0.95, font=\sffamily\footnotesize]
    % Separating hyperplane w^T x + b = 0
    \draw[line width=2pt, color=primarydark] (-3.8, -1.6) -- (3.8, 2.6);
    \node[above, color=primarydark] at (2.4, 2.2) {\textbf{Decision Boundary: } $\mathbf{w}^T \mathbf{x} + b = 0$};

    % Positive margin hyperplane w^T x + b = +1
    \draw[dashed, line width=1.4pt, color=tealborder] (-4.2, -0.5) -- (3.2, 3.7);
    \node[above, color=tealborder] at (-2.5, 1.2) {$\mathbf{w}^T \mathbf{x} + b = +1$};

    % Negative margin hyperplane w^T x + b = -1
    \draw[dashed, line width=1.4pt, color=coralborder] (-3.2, -2.7) -- (4.2, 1.5);
    \node[below, color=coralborder] at (2.5, 0.2) {$\mathbf{w}^T \mathbf{x} + b = -1$};

    % Class +1 data points (teal squares)
    \filldraw[color=tealborder] (-2.2, 2.2) rectangle ++(5pt, 5pt);
    \filldraw[color=tealborder] (0.0, 3.2) rectangle ++(5pt, 5pt);
    \filldraw[color=tealborder] (-1.2, 3.4) rectangle ++(5pt, 5pt);
    
    % Positive support vectors on margin
    \coordinate (SV1) at (-0.6, 1.6);
    \filldraw[color=tealborder] (SV1) rectangle ++(6pt, 6pt);
    \draw[line width=1.5pt, color=forest] ($(SV1) + (3pt, 3pt)$) circle (9pt);
    \node[above, color=forest] at ($(SV1) + (0, 0.35)$) {\textbf{Support Vector } ($\lambda_1^* > 0$)};

    % Class -1 data points (coral circles)
    \filldraw[color=coralborder] (0.6, -2.2) circle (3pt);
    \filldraw[color=coralborder] (2.2, -1.8) circle (3pt);
    \filldraw[color=coralborder] (-1.2, -3.1) circle (3pt);
    
    % Negative support vectors on margin
    \coordinate (SV2) at (0.6, -0.6);
    \filldraw[color=coralborder] (SV2) circle (3.5pt);
    \draw[line width=1.5pt, color=traprad] (SV2) circle (9pt);
    \node[below right, color=traprad] at ($(SV2) + (0.1, -0.1)$) {\textbf{Support Vector } ($\lambda_2^* > 0$)};

    % Margin width arrow
    \draw[<->, >=stealth, line width=1.6pt, color=traprad] (-1.1, 1.3) -- (-0.1, -0.2) 
        node[midway, left] {\textbf{Margin } $\gamma = \frac{2}{\|\mathbf{w}\|_2}$};

    % Normal vector w
    \draw[->, >=stealth, line width=2pt, color=primarydark] (0, 0.5) -- (-0.8, 1.9) 
        node[left] {$\mathbf{w}^* = \sum \lambda_i^* y_i \mathbf{x}_i$};
\end{tikzpicture}
\caption{The Geometry of Support Vector Machines: Maximizing the geometric margin $\frac{2}{\|\mathbf{w}\|_2}$ forces the weight vector to depend strictly on the boundary support vectors. Training samples lying strictly inside the class territories have $\lambda_i^* = 0$ and exert zero influence on the separating hyperplane.}
\label{fig:svm_geometry}
\end{figure}

\subsection{The Primal-Dual Geometry and Duality Gap}
Figure~\ref{fig:duality_gap} illustrates the geometric relationship between the primal optimal value $p^*$, the dual function $g(\boldsymbol{\lambda})$, the duality gap $p^* - d^*$, and Slater's condition.

\begin{figure}[htbp]
\centering
\begin{tikzpicture}[scale=1.05, font=\sffamily\footnotesize]
    % Axes
    \draw[->, thick, color=primarydark!60] (-2.5, 0) -- (3.5, 0) node[right] {$u = f_1(\mathbf{w})$};
    \draw[->, thick, color=primarydark!60] (0, -1.5) -- (0, 4.0) node[above] {$t = f_0(\mathbf{w})$};

    % Achievable set G = {(f1(w), f0(w))}
    \draw[thick, fill=tealborder!15, color=tealborder] 
        (-2.0, 3.5) to[out=-40, in=170] (0.0, 1.8)
                    to[out=-10, in=190] (2.5, 1.2)
                    to[out=60, in=-60] (3.0, 3.8)
                    to[out=170, in=10] cycle;
    \node[color=tealborder!90!black] at (1.5, 2.8) {\large $\mathcal{G} = \{(f_1(\mathbf{w}), f_0(\mathbf{w}))\}$};

    % Primal optimal point p* at u = 0
    \coordinate (Pstar) at (0.0, 1.8);
    \filldraw[color=forest] (Pstar) circle (3pt) node[left=2pt] {$p^* = \inf f_0(\mathbf{w})$};

    % Supporting Hyperplane with slope -lambda
    % Equation: t + lambda * u >= d* => t = d* - lambda * u
    % For lambda = 0.5, d* = 1.0 => t = 1.0 - 0.5 * u
    \draw[very thick, color=coralborder, domain=-2.2:3.2] 
        plot (\x, {1.0 - 0.5*\x}) node[right] {\textbf{Dual Supporting Line } $t + \lambda u = d^*$};

    % Dual intercept d* on vertical axis
    \coordinate (Dstar) at (0, 1.0);
    \filldraw[color=coralborder] (Dstar) circle (3pt) node[below left=2pt] {$d^* = \sup g(\lambda)$};

    % Duality Gap annotation
    \draw[<->, >=stealth, line width=1.5pt, color=traprad] (0.3, 1.0) -- (0.3, 1.8);
    \node[right, color=traprad, align=left] at (0.35, 1.4) {
        \textbf{Duality Gap } $p^* - d^* \ge 0$ \\
        \scriptsize (Zero if Slater's Condition holds)
    };
\end{tikzpicture}
\caption{The Geometry of Lagrange Duality: The set $\mathcal{G}$ represents the set of achievable constraint-objective values $(f_1(\mathbf{w}), f_0(\mathbf{w}))$. The primal minimum $p^*$ is the lowest point of $\mathcal{G}$ on or to the left of the vertical axis ($u \le 0$). The dual function $g(\lambda)$ represents the intercept of a supporting hyperplane with normal $(\lambda, 1)$ supporting $\mathcal{G}$ from below. Weak duality states that $d^* \le p^*$. Under Slater's condition, the supporting hyperplane touches $\mathcal{G}$ directly at $p^*$, collapsing the duality gap to zero ($p^* = d^*$).}
\label{fig:duality_gap}
\end{figure}

\subsection{Complementary Slackness as a Geometric Switch}
Figure~\ref{fig:slackness_switch} reveals how complementary slackness partitions parameter space into inactive non-support vectors and active margin constraints.

\begin{figure}[htbp]
\centering
\begin{tikzpicture}[scale=0.95, font=\sffamily\footnotesize]
    % Panel 1: Inactive Constraint (Non-Support Vector)
    \begin{scope}[shift={(-4.5, 0)}]
        \draw[thick, fill=primarydark!10, color=primarydark] (-2.2, -2.0) rectangle (2.2, 2.0);
        \node[above, color=primarydark] at (0, 2.0) {\textbf{Feasible Region } $f_i(\mathbf{w}) \le 0$};

        % Solution w* inside interior
        \coordinate (W1) at (0, 0.2);
        \filldraw[color=forest] (W1) circle (3pt) node[above] {$\mathbf{w}^*$};
        \node[below, color=forest, align=center] at (0, -0.2) {$f_i(\mathbf{w}^*) < 0$ \\ (\textbf{Strictly Interior})};

        % Constraint boundary far away
        \draw[line width=1.8pt, color=traprad] (1.8, -2.0) -- (1.8, 2.0);
        \node[rotate=90, above, color=traprad] at (1.8, 0) {\scriptsize Boundary $f_i = 0$};

        \node[draw=forest, fill=white, rounded corners=3pt, align=center] at (0, -2.8) {
            \textbf{Inactive Constraint:} \\
            $\lambda_i^* = 0$ \\
            \scriptsize Zero contact force
        };
    \end{scope}

    % Panel 2: Active Constraint (Support Vector)
    \begin{scope}[shift={(4.5, 0)}]
        \draw[thick, fill=primarydark!10, color=primarydark] (-2.2, -2.0) rectangle (2.2, 2.0);
        \node[above, color=primarydark] at (0, 2.0) {\textbf{Feasible Region } $f_i(\mathbf{w}) \le 0$};

        % Solution w* directly on boundary
        \coordinate (W2) at (1.8, 0.2);
        \filldraw[color=traprad] (W2) circle (3pt);
        \node[left, color=traprad, align=right] at (1.7, 0.2) {$\mathbf{w}^*$ on boundary \\ $f_i(\mathbf{w}^*) = 0$};

        % Boundary
        \draw[line width=1.8pt, color=traprad] (1.8, -2.0) -- (1.8, 2.0);
        \node[rotate=90, above, color=traprad] at (1.8, 1.2) {\scriptsize Boundary $f_i = 0$};

        % Normal force vector lambda * grad f
        \draw[->, >=stealth, line width=2pt, color=coralborder] (W2) -- ($(W2) + (-1.4, 0)$) 
            node[below left] {$\lambda_i^* \nabla f_i$};

        \node[draw=traprad, fill=white, rounded corners=3pt, align=center] at (0, -2.8) {
            \textbf{Active Constraint:} \\
            $\lambda_i^* > 0$ \\
            \scriptsize Support Vector reaction force
        };
    \end{scope}
\end{tikzpicture}
\caption{The Complementary Slackness Geometric Switch: For every inequality constraint, $\lambda_i^* f_i(\mathbf{w}^*) = 0$. Left: If the optimal parameter vector $\mathbf{w}^*$ lies strictly in the interior ($f_i < 0$), the constraint is non-binding and the dual multiplier is identically zero ($\lambda_i^* = 0$). Right: If the multiplier is strictly positive ($\lambda_i^* > 0$), the parameter vector is pressed firmly against the boundary ($f_i = 0$), defining an active support vector.}
\label{fig:slackness_switch}
\end{figure}

% ----------------------------------------------------------------------
% PART 4: DEEP INTUITION & CONCEPTUAL DERIVATIONS
% ----------------------------------------------------------------------
\section{Deep Intuition \& Conceptual Derivations}

\subsection{Derivation of the Dual Support Vector Machine}
Given a training dataset $\mathcal{D} = \{(\mathbf{x}_i, y_i)\}_{i=1}^N$ with feature vectors $\mathbf{x}_i \in \mathbb{R}^d$ and binary class labels $y_i \in \{-1, +1\}$, a linear separating hyperplane is defined by normal vector $\mathbf{w} \in \mathbb{R}^d$ and bias intercept $b \in \mathbb{R}$:
\begin{equation}
\mathcal{H} = \{\mathbf{x} \in \mathbb{R}^d \mid \mathbf{w}^T \mathbf{x} + b = 0\}
\end{equation}
The signed perpendicular Euclidean distance from any sample $\mathbf{x}_i$ to the hyperplane is:
\begin{equation}
\gamma_i = \frac{y_i (\mathbf{w}^T \mathbf{x}_i + b)}{\|\mathbf{w}\|_2}
\end{equation}
The \textbf{canonical hyperplane} fixes the scale of $(\mathbf{w}, b)$ such that the closest training samples to the boundary achieve a functional margin of exactly 1:
\begin{equation}
\min_{i=1,\dots,N} y_i (\mathbf{w}^T \mathbf{x}_i + b) = 1
\end{equation}
Then the total geometric margin separating the two classes is:
\begin{equation}
\gamma = \frac{1}{\|\mathbf{w}\|_2} + \frac{1}{\|\mathbf{w}\|_2} = \frac{2}{\|\mathbf{w}\|_2}
\end{equation}
Maximizing the geometric margin $\frac{2}{\|\mathbf{w}\|_2}$ is mathematically equivalent to minimizing $\frac{1}{2}\|\mathbf{w}\|_2^2$:
\begin{equation}
\min_{\mathbf{w} \in \mathbb{R}^d, b \in \mathbb{R}} \frac{1}{2} \|\mathbf{w}\|_2^2 \quad \text{subject to } y_i (\mathbf{w}^T \mathbf{x}_i + b) \ge 1 \quad (i = 1, \dots, N)
\label{eq:hard_svm_primal}
\end{equation}

\subsubsection{Formulating the Lagrangian}
Rewriting the $N$ margin constraints in standard format $1 - y_i(\mathbf{w}^T \mathbf{x}_i + b) \le 0$:
\begin{equation}
\mathcal{L}(\mathbf{w}, b, \boldsymbol{\lambda}) = \frac{1}{2}\|\mathbf{w}\|_2^2 - \sum_{i=1}^N \lambda_i \left[ y_i(\mathbf{w}^T \mathbf{x}_i + b) - 1 \right]
\label{eq:svm_lagrangian}
\end{equation}
with dual multipliers $\lambda_i \ge 0$.

\subsubsection{Primal Stationarity Conditions}
Setting the partial gradients of $\mathcal{L}$ with respect to the primal parameters $\mathbf{w}$ and $b$ to zero:
\begin{align}
\nabla_{\mathbf{w}} \mathcal{L} &= \mathbf{w} - \sum_{i=1}^N \lambda_i y_i \mathbf{x}_i = \mathbf{0} \implies \mathbf{w}^* = \sum_{i=1}^N \lambda_i y_i \mathbf{x}_i \label{eq:svm_w_stat} \\
\frac{\partial \mathcal{L}}{\partial b} &= -\sum_{i=1}^N \lambda_i y_i = 0 \implies \sum_{i=1}^N \lambda_i y_i = 0 \label{eq:svm_b_stat}
\end{align}

\subsubsection{Constructing the Dual Objective}
Substituting $\mathbf{w} = \sum_{i=1}^N \lambda_i y_i \mathbf{x}_i$ and $\sum_{i=1}^N \lambda_i y_i = 0$ back into the Lagrangian:
\begin{align}
\mathcal{L} &= \frac{1}{2} \left( \sum_{i=1}^N \lambda_i y_i \mathbf{x}_i \right)^T \left( \sum_{j=1}^N \lambda_j y_j \mathbf{x}_j \right) - \sum_{i=1}^N \lambda_i y_i \left( \sum_{j=1}^N \lambda_j y_j \mathbf{x}_j \right)^T \mathbf{x}_i - b \underbrace{\sum_{i=1}^N \lambda_i y_i}_{= 0} + \sum_{i=1}^N \lambda_i \nonumber \\
&= \frac{1}{2} \sum_{i=1}^N \sum_{j=1}^N \lambda_i \lambda_j y_i y_j (\mathbf{x}_i^T \mathbf{x}_j) - \sum_{i=1}^N \sum_{j=1}^N \lambda_i \lambda_j y_i y_j (\mathbf{x}_i^T \mathbf{x}_j) + \sum_{i=1}^N \lambda_i \nonumber \\
&= \sum_{i=1}^N \lambda_i - \frac{1}{2} \sum_{i=1}^N \sum_{j=1}^N \lambda_i \lambda_j y_i y_j (\mathbf{x}_i^T \mathbf{x}_j)
\end{align}
This yields the celebrated \textbf{Dual Support Vector Machine Quadratic Program}:
\begin{equation}
\max_{\boldsymbol{\lambda} \in \mathbb{R}^N} \sum_{i=1}^N \lambda_i - \frac{1}{2} \sum_{i=1}^N \sum_{j=1}^N \lambda_i \lambda_j y_i y_j (\mathbf{x}_i^T \mathbf{x}_j) \quad \text{subject to } \begin{cases} \lambda_i \ge 0, & \forall i \in \{1, \dots, N\} \\ \sum_{i=1}^N \lambda_i y_i = 0 & \end{cases}
\label{eq:svm_dual_qp}
\end{equation}

\subsection{Soft-Margin SVM and the Box Constraint $0 \le \lambda_i \le C$}
When classes overlap in feature space, no linear separating hyperplane exists. The Soft-Margin SVM introduces non-negative slack variables $\xi_i \ge 0$:
\begin{equation}
\min_{\mathbf{w}, b, \boldsymbol{\xi}} \frac{1}{2}\|\mathbf{w}\|_2^2 + C \sum_{i=1}^N \xi_i \quad \text{subject to } \begin{cases} y_i(\mathbf{w}^T \mathbf{x}_i + b) \ge 1 - \xi_i, & \forall i \\ \xi_i \ge 0, & \forall i \end{cases}
\label{eq:soft_svm_primal}
\end{equation}
The Lagrangian introduces multipliers $\lambda_i \ge 0$ for the margin constraint and $\mu_i \ge 0$ for the non-negativity of slack:
\begin{equation}
\mathcal{L}(\mathbf{w}, b, \boldsymbol{\xi}, \boldsymbol{\lambda}, \boldsymbol{\mu}) = \frac{1}{2}\|\mathbf{w}\|_2^2 + C \sum_{i=1}^N \xi_i - \sum_{i=1}^N \lambda_i \left[ y_i(\mathbf{w}^T \mathbf{x}_i + b) - 1 + \xi_i \right] - \sum_{i=1}^N \mu_i \xi_i
\end{equation}
Stationarity with respect to the slack variable $\xi_i$:
\begin{equation}
\frac{\partial \mathcal{L}}{\partial \xi_i} = C - \lambda_i - \mu_i = 0 \implies \lambda_i + \mu_i = C \implies \lambda_i = C - \mu_i
\end{equation}
Since $\mu_i \ge 0$, it follows that $\lambda_i \le C$. Combined with $\lambda_i \ge 0$, this establishes the famous \textbf{box constraint}:
\begin{equation}
0 \le \lambda_i \le C \quad \forall i \in \{1, \dots, N\}
\label{eq:box_constraint}
\end{equation}
The dual objective function for Soft-Margin SVM is identical to the Hard-Margin dual~\eqref{eq:svm_dual_qp}; the only mathematical modification is that the upper boundary of the multipliers is clipped at $C$.

\begin{ponderbox}{The Soft-Margin Trade-Off: What Does $C$ Actually Control?}
What is the fundamental bias-variance trade-off governed by the SVM penalty hyperparameter $C$?
\tcblower
\textbf{The Socratic Resolution:} Hyperparameter $C$ acts as the \textbf{economic exchange rate} between margin width and empirical training error:
\begin{itemize}
    \item \textbf{Large $C$ ($C \to \infty$):} Massive penalty on slack violations. The algorithm prioritizes classifying every training point correctly, forcing slacks $\xi_i \to 0$. The margin width narrows, making the model highly sensitive to individual outliers and risking severe overfitting.
    \item \textbf{Small $C$ ($C \to 0$):} Minimal penalty on slack violations. The algorithm prioritizes maximizing the margin width $\frac{2}{\|\mathbf{w}\|_2}$, readily tolerating training errors and margin violations to find a broader, more robust separating corridor. This increases regularization and prevents overfitting.
\end{itemize}
\end{ponderbox}

\begin{ponderbox}{Primal vs.\ Dual Computational Complexity: When to Solve Which?}
When should an engineer solve the Primal SVM formulation, and when should they solve the Dual SVM formulation?
\tcblower
\textbf{The Socratic Resolution:} The choice depends strictly on the ratio between the number of training samples $N$ and the feature dimensionality $d$:
\begin{itemize}
    \item \textbf{Primal SVM ($\mathcal{O}(d)$ to $\mathcal{O}(d^3)$):} The primal variable vector $\mathbf{w} \in \mathbb{R}^d$ has dimension $d$. When training on massive tabular or text datasets with millions of samples ($N \gg d$, e.g., $N = 10^7, d = 100$), solving the primal problem via subgradient descent or stochastic dual coordinate ascent takes linear time $\mathcal{O}(N d)$ without storing an $N \times N$ matrix.
    \item \textbf{Dual SVM ($\mathcal{O}(N^2)$ to $\mathcal{O}(N^3)$):} The dual variable vector $\boldsymbol{\lambda} \in \mathbb{R}^N$ has dimension $N$. When the feature space is very high-dimensional or infinite ($d \gg N$, e.g., genomics, RBF kernels where $d = \infty$), the primal cannot be solved. The dual depends strictly on the $N \times N$ Gram kernel matrix $K_{ij} = k(\mathbf{x}_i, \mathbf{x}_j)$, making non-linear classification computationally feasible.
\end{itemize}
\end{ponderbox}

% ----------------------------------------------------------------------
% PART 5: MULTI-PAGE WORKED EXAMPLES
% ----------------------------------------------------------------------
\section{Multi-Page Breakable Numerical Worked Examples}

\begin{examplebox}[breakable]{Step-by-Step Resolution of a Constrained Quadratic via KKT Conditions}
\textbf{Problem Statement:}
Consider the continuous optimization problem:
\begin{equation}
\min_{x, y \in \mathbb{R}} f_0(x, y) = x^2 + y^2 \quad \text{subject to } x + y \ge 4
\end{equation}
Rewrite the inequality constraint in standard format $f_1(x, y) = 4 - x - y \le 0$.

\textbf{Task Requirements:}
\begin{enumerate}
    \item Formulate the Lagrangian $\mathcal{L}(x, y, \lambda)$ with dual variable $\lambda \ge 0$.
    \item Explicitly write down the four KKT conditions.
    \item Test Case 1: $\lambda = 0$ (inactive constraint). Determine the solution and prove why it violates primal feasibility.
    \item Test Case 2: $\lambda > 0$ (active constraint). Analytically compute $(x^*, y^*)$ and $\lambda^*$, and verify that all four KKT conditions are rigorously satisfied.
    \item Compute the optimal primal value $p^*$ and dual value $d^*$, and verify that the duality gap is zero.
\end{enumerate}

\tcblower

\textbf{Step 1: The Lagrangian}
\begin{equation}
\mathcal{L}(x, y, \lambda) = x^2 + y^2 + \lambda(4 - x - y)
\end{equation}
with dual feasibility requirement $\lambda \ge 0$.

\textbf{Step 2: The Four KKT Conditions}
\begin{enumerate}
    \item \textbf{Primal Feasibility:}
    \begin{equation}
    4 - x - y \le 0 \iff x + y \ge 4
    \label{eq:ex_kkt_primal}
    \end{equation}
    \item \textbf{Dual Feasibility:}
    \begin{equation}
    \lambda \ge 0
    \label{eq:ex_kkt_dual}
    \end{equation}
    \item \textbf{Complementary Slackness:}
    \begin{equation}
    \lambda(4 - x - y) = 0
    \label{eq:ex_kkt_slack}
    \end{equation}
    \item \textbf{Stationarity:}
    \begin{align}
    \frac{\partial \mathcal{L}}{\partial x} &= 2x - \lambda = 0 \implies x = \frac{\lambda}{2} \label{eq:ex_kkt_stat_x} \\
    \frac{\partial \mathcal{L}}{\partial y} &= 2y - \lambda = 0 \implies y = \frac{\lambda}{2} \label{eq:ex_kkt_stat_y}
    \end{align}
\end{enumerate}

\textbf{Step 3: Hypothesis Testing -- Case 1: $\lambda = 0$ (Inactive Constraint)}
Assume $\lambda = 0$.
From the stationarity conditions~\eqref{eq:ex_kkt_stat_x} and~\eqref{eq:ex_kkt_stat_y}:
\begin{equation}
x = \frac{0}{2} = 0, \quad y = \frac{0}{2} = 0
\end{equation}
Now check Primal Feasibility~\eqref{eq:ex_kkt_primal}:
\begin{equation}
4 - x - y = 4 - 0 - 0 = 4 \le 0 \quad \text{\textbf{(CONTRADICTION! False)}}
\end{equation}
The point $(0, 0)$ lies in the infeasible halfplane. Hence $\lambda = 0$ is rejected.

\textbf{Step 4: Hypothesis Testing -- Case 2: $\lambda > 0$ (Active Constraint)}
Because $\lambda > 0$, Complementary Slackness~\eqref{eq:ex_kkt_slack} strictly forces the constraint boundary to be active:
\begin{equation}
4 - x - y = 0 \implies x + y = 4
\end{equation}
Substitute the stationarity relations $x = \frac{\lambda}{2}$ and $y = \frac{\lambda}{2}$:
\begin{equation}
\frac{\lambda}{2} + \frac{\lambda}{2} = 4 \implies \lambda = 4
\end{equation}
Back-substituting $\lambda^* = 4$:
\begin{align}
x^* &= \frac{4}{2} = 2.0 \\
y^* &= \frac{4}{2} = 2.0
\end{align}

\textbf{Step 5: Full KKT Verification}
\begin{itemize}
    \item Primal Feasibility: $4 - x^* - y^* = 4 - 2 - 2 = 0 \le 0$ (\textbf{Satisfied}).
    \item Dual Feasibility: $\lambda^* = 4 \ge 0$ (\textbf{Satisfied}).
    \item Complementary Slackness: $\lambda^*(4 - x^* - y^*) = 4(0) = 0$ (\textbf{Satisfied}).
    \item Stationarity: $\nabla f_0(2, 2) + \lambda^* \nabla f_1(2, 2) = [4, 4]^T + 4[-1, -1]^T = [0, 0]^T$ (\textbf{Satisfied}).
\end{itemize}
Because the objective $f_0(x, y) = x^2 + y^2$ is strictly convex and the constraint $f_1(x, y) = 4 - x - y$ is convex (affine), the KKT conditions are both necessary and sufficient.
The unique global minimizer is $(x^*, y^*) = (2, 2)$.

\textbf{Step 6: Optimal Objective and Duality Gap}
The primal optimal value is:
\begin{equation}
p^* = f_0(2, 2) = 2^2 + 2^2 = 4 + 4 = 8.000
\end{equation}
The Lagrange dual function is:
\begin{align}
g(\lambda) &= \inf_{x, y} \left( x^2 + y^2 + \lambda(4 - x - y) \right) \nonumber \\
&= \left(\frac{\lambda}{2}\right)^2 + \left(\frac{\lambda}{2}\right)^2 + \lambda\left(4 - \frac{\lambda}{2} - \frac{\lambda}{2}\right) \nonumber \\
&= \frac{\lambda^2}{4} + \frac{\lambda^2}{4} + 4\lambda - \lambda^2 = 4\lambda - \frac{\lambda^2}{2}
\end{align}
Maximizing the dual function over $\lambda \ge 0$:
\begin{equation}
g'(\lambda) = 4 - \lambda = 0 \implies \lambda^* = 4
\end{equation}
The optimal dual value is:
\begin{equation}
d^* = g(4) = 4(4) - \frac{4^2}{2} = 16 - 8 = 8.000
\end{equation}
The duality gap is:
\begin{equation}
p^* - d^* = 8.000 - 8.000 = 0.000
\end{equation}
Strong duality holds with zero duality gap.
\end{examplebox}

\begin{examplebox}[breakable]{Analytical 3-Point Hard-Margin Support Vector Machine}
\textbf{Problem Statement:}
Consider a two-dimensional classification dataset comprising $N = 3$ points:
\begin{itemize}
    \item $\mathbf{x}_1 = [1, 1]^T$ with class label $y_1 = +1$
    \item $\mathbf{x}_2 = [2, 0]^T$ with class label $y_2 = +1$
    \item $\mathbf{x}_3 = [0, 0]^T$ with class label $y_3 = -1$
\end{itemize}

\textbf{Task Requirements:}
\begin{enumerate}
    \item Compute all pairwise inner products $\mathbf{x}_i^T \mathbf{x}_j$ and formulate the Gram kernel matrix $K$.
    \item Enforce the equality constraint $\sum_{i=1}^3 \lambda_i y_i = 0$ to eliminate $\lambda_3$.
    \item Express the dual objective as a bivariate quadratic function $D(\lambda_1, \lambda_2)$ and compute its unconstrained maximum.
    \item Determine which training samples are support vectors.
    \item Reconstruct the optimal weight vector $\mathbf{w}^*$, compute the bias intercept $b^*$, and write the equation of the decision boundary.
\end{enumerate}

\tcblower

\textbf{Step 1: Pairwise Inner Products and Gram Matrix}
\begin{align}
\mathbf{x}_1^T \mathbf{x}_1 &= 1(1) + 1(1) = 2 \\
\mathbf{x}_2^T \mathbf{x}_2 &= 2(2) + 0(0) = 4 \\
\mathbf{x}_3^T \mathbf{x}_3 &= 0(0) + 0(0) = 0 \\
\mathbf{x}_1^T \mathbf{x}_2 &= 1(2) + 1(0) = 2 \\
\mathbf{x}_1^T \mathbf{x}_3 &= 1(0) + 1(0) = 0 \\
\mathbf{x}_2^T \mathbf{x}_3 &= 2(0) + 0(0) = 0
\end{align}

\textbf{Step 2: Elimination of $\lambda_3$}
The equality constraint requires:
\begin{equation}
\sum_{i=1}^3 \lambda_i y_i = \lambda_1(+1) + \lambda_2(+1) + \lambda_3(-1) = 0 \implies \lambda_3 = \lambda_1 + \lambda_2
\end{equation}

\textbf{Step 3: Formulation of the Dual Objective}
The dual objective function is:
\begin{equation}
D(\boldsymbol{\lambda}) = \sum_{i=1}^3 \lambda_i - \frac{1}{2}\sum_{i=1}^3 \sum_{j=1}^3 \lambda_i \lambda_j y_i y_j (\mathbf{x}_i^T \mathbf{x}_j)
\end{equation}
Substitute $\lambda_3 = \lambda_1 + \lambda_2$:
\begin{itemize}
    \item Linear term:
    \begin{equation}
    \lambda_1 + \lambda_2 + \lambda_3 = \lambda_1 + \lambda_2 + (\lambda_1 + \lambda_2) = 2(\lambda_1 + \lambda_2)
    \end{equation}
    \item Quadratic term:
    \begin{align}
    \frac{1}{2} \sum_{i,j} \lambda_i \lambda_j y_i y_j (\mathbf{x}_i^T \mathbf{x}_j) &= \frac{1}{2} \left[ \lambda_1^2 (1)^2 (2) + \lambda_2^2 (1)^2 (4) + \lambda_3^2 (-1)^2 (0) + 2\lambda_1 \lambda_2 (1)(1)(2) - 0 - 0 \right] \nonumber \\
    &= \frac{1}{2} \left[ 2\lambda_1^2 + 4\lambda_2^2 + 4\lambda_1 \lambda_2 \right] = \lambda_1^2 + 2\lambda_2^2 + 2\lambda_1 \lambda_2
    \end{align}
\end{itemize}
Thus:
\begin{equation}
D(\lambda_1, \lambda_2) = 2\lambda_1 + 2\lambda_2 - \lambda_1^2 - 2\lambda_2^2 - 2\lambda_1 \lambda_2
\end{equation}

\textbf{Step 4: Maximization with Respect to $\lambda_1, \lambda_2$}
Setting partial derivatives to zero:
\begin{align}
\frac{\partial D}{\partial \lambda_1} &= 2 - 2\lambda_1 - 2\lambda_2 = 0 \implies \lambda_1 + \lambda_2 = 1 \label{eq:svm_stat1} \\
\frac{\partial D}{\partial \lambda_2} &= 2 - 4\lambda_2 - 2\lambda_1 = 0 \implies \lambda_1 + 2\lambda_2 = 1 \label{eq:svm_stat2}
\end{align}
Subtracting~\eqref{eq:svm_stat1} from~\eqref{eq:svm_stat2}:
\begin{equation}
(\lambda_1 + 2\lambda_2) - (\lambda_1 + \lambda_2) = 1 - 1 \implies \lambda_2^* = 0
\end{equation}
Substituting $\lambda_2^* = 0$ into~\eqref{eq:svm_stat1}:
\begin{equation}
\lambda_1^* = 1.0
\end{equation}
Finally, recovering $\lambda_3^*$:
\begin{equation}
\lambda_3^* = \lambda_1^* + \lambda_2^* = 1.0 + 0 = 1.0
\end{equation}
Notice that all dual constraints are satisfied: $\lambda_1^* = 1 \ge 0$, $\lambda_2^* = 0 \ge 0$, and $\lambda_3^* = 1 \ge 0$.

\textbf{Step 5: Support Vectors and Boundary Reconstruction}
\begin{itemize}
    \item $\lambda_1^* = 1 > 0 \implies \mathbf{x}_1 = [1, 1]^T$ is an active \textbf{Support Vector}.
    \item $\lambda_2^* = 0 \implies \mathbf{x}_2 = [2, 0]^T$ is a \textbf{Non-Support Vector} (interior point).
    \item $\lambda_3^* = 1 > 0 \implies \mathbf{x}_3 = [0, 0]^T$ is an active \textbf{Support Vector}.
\end{itemize}
Now reconstruct the normal weight vector $\mathbf{w}^*$:
\begin{align}
\mathbf{w}^* &= \sum_{i=1}^3 \lambda_i^* y_i \mathbf{x}_i \nonumber \\
&= (1)(+1)\begin{bmatrix} 1 \\ 1 \end{bmatrix} + (0)(+1)\begin{bmatrix} 2 \\ 0 \end{bmatrix} + (1)(-1)\begin{bmatrix} 0 \\ 0 \end{bmatrix} \nonumber \\
&= \begin{bmatrix} 1 \\ 1 \end{bmatrix} + \begin{bmatrix} 0 \\ 0 \end{bmatrix} + \begin{bmatrix} 0 \\ 0 \end{bmatrix} = \begin{bmatrix} 1.0 \\ 1.0 \end{bmatrix}
\end{align}
Compute the bias intercept $b^*$ using active support vector $\mathbf{x}_1$:
\begin{equation}
y_1(\mathbf{w}^{*T} \mathbf{x}_1 + b^*) = 1 \implies 1\left( \begin{bmatrix} 1 & 1 \end{bmatrix} \begin{bmatrix} 1 \\ 1 \end{bmatrix} + b^* \right) = 1 \implies 2 + b^* = 1 \implies b^* = -1.0
\end{equation}
Verify on negative support vector $\mathbf{x}_3$:
\begin{equation}
y_3(\mathbf{w}^{*T} \mathbf{x}_3 + b^*) = -1\left( \begin{bmatrix} 1 & 1 \end{bmatrix} \begin{bmatrix} 0 \\ 0 \end{bmatrix} - 1 \right) = -1(0 - 1) = 1 \quad (\text{Exact margin})
\end{equation}
The optimal separating hyperplane equation is:
\begin{equation}
x_1 + x_2 - 1 = 0
\end{equation}
The geometric margin width is $\gamma = \frac{2}{\|\mathbf{w}^*\|_2} = \frac{2}{\sqrt{1^2 + 1^2}} = \frac{2}{\sqrt{2}} = \sqrt{2} \approx 1.414$.
\end{examplebox}

\begin{examplebox}[breakable]{Step-by-Step Soft-Margin SVM with Margin Violator and Slack Analysis}
\textbf{Problem Statement:}
Consider a one-dimensional dataset with $N = 4$ points and penalty parameter $C = 2.0$:
\begin{itemize}
    \item $x_1 = 2.0, y_1 = +1$ (Positive class)
    \item $x_2 = 4.0, y_2 = +1$ (Positive class, deep interior)
    \item $x_3 = -2.0, y_3 = -1$ (Negative class)
    \item $x_4 = 1.0, y_4 = -1$ (Outlier: negative label, but located on the positive axis!)
\end{itemize}

\textbf{Task Requirements:}
\begin{enumerate}
    \item Formulate the dual constraints on $\lambda_1, \lambda_2, \lambda_3, \lambda_4$ incorporating the box constraint $\lambda_i \in [0, 2.0]$.
    \item Identify the status of outlier $x_4$ and explain why it forces $\lambda_4^* = C = 2.0$.
    \item Using complementary slackness, deduce which points are free support vectors, non-support vectors, and bounded support vectors.
    \item Solve for the exact dual multipliers $\boldsymbol{\lambda}^*$, optimal weight $w^*$, and bias intercept $b^*$.
    \item Calculate the exact slack variables $\xi_1^*, \xi_2^*, \xi_3^*, \xi_4^*$ and verify all KKT conditions.
\end{enumerate}

\tcblower

\textbf{Step 1: Dual Constraints and Outlier Identification}
The dual box constraints are:
\begin{equation}
0 \le \lambda_i \le C = 2.0 \quad \forall i \in \{1, 2, 3, 4\}
\end{equation}
The equality constraint is:
\begin{equation}
\sum_{i=1}^4 \lambda_i y_i = \lambda_1(+1) + \lambda_2(+1) + \lambda_3(-1) + \lambda_4(-1) = 0 \implies \lambda_1 + \lambda_2 = \lambda_3 + \lambda_4
\end{equation}
Sample $x_4 = 1.0$ has negative label $y_4 = -1$, but is located to the right of $x_3 = -2.0$ and between $x_3$ and $x_1 = 2.0$. Because a single threshold cannot separate $x_4$ without severe margin violation, $x_4$ is an outlier. For any non-trivial margin, $x_4$ violates the margin boundary ($y_4(w x_4 + b) < 1$), which demands slack $\xi_4 > 0$.
By complementary slackness $(C - \lambda_4)\xi_4 = 0$, having $\xi_4 > 0$ strictly forces:
\begin{equation}
\lambda_4^* = C = 2.0
\end{equation}

\textbf{Step 2: Regime Classification of Remaining Points}
\begin{itemize}
    \item $x_2 = 4.0$ lies far in the positive interior. It is a non-support vector: $\lambda_2^* = 0 \implies \xi_2^* = 0$.
    \item Substituting $\lambda_2 = 0$ and $\lambda_4 = 2.0$ into the equality constraint:
    \begin{equation}
    \lambda_1 + 0 = \lambda_3 + 2.0 \implies \lambda_1 = \lambda_3 + 2.0
    \end{equation}
    Because $\lambda_1 \le C = 2.0$ and $\lambda_3 \ge 0$, this requires:
    \begin{equation}
    \lambda_3^* = 0, \quad \lambda_1^* = 2.0
    \end{equation}
    Wait! If $\lambda_3^* = 0$, then $x_3 = -2.0$ would exert zero weight, leaving no negative support vectors on the negative margin. Let us solve the dual objective analytically.
\end{itemize}

\textbf{Step 3: Analytical Dual Maximization}
With $\lambda_2 = 0$, the equality constraint is $\lambda_1 = \lambda_3 + \lambda_4$.
Kernel products in 1D ($x_i x_j$):
\begin{itemize}
    \item $x_1^2 = 4$, $x_3^2 = 4$, $x_4^2 = 1$
    \item $x_1 x_3 = -4$, $x_1 x_4 = 2$, $x_3 x_4 = -2$
\end{itemize}
The dual quadratic objective is:
\begin{align}
D(\boldsymbol{\lambda}) &= \sum \lambda_i - \frac{1}{2}\sum_{i,j} \lambda_i \lambda_j y_i y_j x_i x_j \nonumber \\
&= \sum \lambda_i - \frac{1}{2}\left( \sum \lambda_i y_i x_i \right)^2 \nonumber \\
&= (\lambda_1 + \lambda_3 + \lambda_4) - \frac{1}{2} w^2
\end{align}
where $w = \lambda_1(1)(2) + \lambda_3(-1)(-2) + \lambda_4(-1)(1) = 2\lambda_1 + 2\lambda_3 - \lambda_4$.
Substituting $\lambda_1 = \lambda_3 + \lambda_4$:
\begin{equation}
w = 2(\lambda_3 + \lambda_4) + 2\lambda_3 - \lambda_4 = 4\lambda_3 + \lambda_4
\end{equation}
Linear term:
\begin{equation}
\lambda_1 + \lambda_3 + \lambda_4 = (\lambda_3 + \lambda_4) + \lambda_3 + \lambda_4 = 2\lambda_3 + 2\lambda_4
\end{equation}
The dual function becomes:
\begin{equation}
D(\lambda_3, \lambda_4) = 2\lambda_3 + 2\lambda_4 - \frac{1}{2}(4\lambda_3 + \lambda_4)^2 = 2\lambda_3 + 2\lambda_4 - 8\lambda_3^2 - \frac{1}{2}\lambda_4^2 - 4\lambda_3 \lambda_4
\end{equation}
Differentiating with respect to $\lambda_3$:
\begin{equation}
\frac{\partial D}{\partial \lambda_3} = 2 - 16\lambda_3 - 4\lambda_4 = 0 \implies 4\lambda_3 + \lambda_4 = \frac{1}{2}
\end{equation}
Notice that $w = 4\lambda_3 + \lambda_4 = 0.50$.
Now differentiate with respect to $\lambda_4$:
\begin{equation}
\frac{\partial D}{\partial \lambda_4} = 2 - \lambda_4 - 4\lambda_3 = 2 - (4\lambda_3 + \lambda_4) = 2 - \frac{1}{2} = 1.5 > 0
\end{equation}
Because $\frac{\partial D}{\partial \lambda_4} > 0$ everywhere, $D$ is strictly increasing in $\lambda_4$. To maximize $D$, $\lambda_4$ must be pushed to its absolute upper box bound:
\begin{equation}
\lambda_4^* = C = 2.0
\end{equation}
Now re-solve for $\lambda_3$ on the boundary $\lambda_4 = 2.0$:
\begin{equation}
\frac{\partial D}{\partial \lambda_3}\Big|_{\lambda_4 = 2.0} = 2 - 16\lambda_3 - 4(2.0) = 2 - 16\lambda_3 - 8 = -6 - 16\lambda_3 < 0 \quad \forall \lambda_3 \ge 0
\end{equation}
Because the derivative is strictly negative for all $\lambda_3 \ge 0$, the maximum occurs at the lower bound:
\begin{equation}
\lambda_3^* = 0.0
\end{equation}
Then by the equality constraint:
\begin{equation}
\lambda_1^* = \lambda_3^* + \lambda_4^* = 0.0 + 2.0 = 2.0
\end{equation}
Thus, the optimal dual multipliers are:
\begin{equation}
\lambda_1^* = 2.0, \quad \lambda_2^* = 0.0, \quad \lambda_3^* = 0.0, \quad \lambda_4^* = 2.0
\end{equation}

\textbf{Step 4: Primal Parameters and Hyperplane}
Compute the optimal weight $w^*$:
\begin{equation}
w^* = \lambda_1^* y_1 x_1 + \lambda_4^* y_4 x_4 = (2.0)(+1)(2.0) + (2.0)(-1)(1.0) = 4.0 - 2.0 = 2.0
\end{equation}
Now determine the bias intercept $b^*$. 
Since $x_1$ is at the boundary $\lambda_1^* = C$, it satisfies $y_1(w^* x_1 + b^*) = 1 - \xi_1^*$. For balanced separation between the positive cluster and the negative outlier:
Set the decision threshold midway between $x_4 = 1.0$ and $x_1 = 2.0$:
$x_{\text{threshold}} = 1.5 \implies w^*(1.5) + b^* = 0 \implies 2.0(1.5) + b^* = 0 \implies b^* = -3.0$.

\textbf{Step 5: Slack Variable Verification}
Evaluate the functional margins $y_i(w^* x_i + b^*)$:
\begin{itemize}
    \item Sample 1: $y_1(2(2.0) - 3) = +1(4 - 3) = 1.0 \implies \xi_1^* = 0$ (Lies exactly on the margin).
    \item Sample 2: $y_2(2(4.0) - 3) = +1(8 - 3) = 5.0 > 1 \implies \xi_2^* = 0$ (Deep interior, $\lambda_2^* = 0$).
    \item Sample 3: $y_3(2(-2.0) - 3) = -1(-4 - 3) = +7.0 > 1 \implies \xi_3^* = 0$ (Deep interior, $\lambda_3^* = 0$).
    \item Sample 4 (Outlier): $y_4(2(1.0) - 3) = -1(2 - 3) = -1(-1) = +1.0$.
    Functional margin is $1.0 = 1 - \xi_4 \implies \xi_4^* = 0$.
\end{itemize}
All constraints and complementary slackness conditions are rigorously satisfied.
\end{examplebox}

% ----------------------------------------------------------------------
% PART 6: EXPLICIT GOTCHAS & DIAGNOSTIC MISCONCEPTIONS
% ----------------------------------------------------------------------
\section{Explicit Gotchas \& Diagnostic Misconceptions}

\begin{trapbox}{Gotcha 1: Inverting Multiplier Signs for Inequality Constraints}
A student formulates the Lagrangian for inequality constraint $f_i(\mathbf{w}) \le 0$ as $\mathcal{L} = f_0 - \sum \lambda_i f_i$ with $\lambda_i \ge 0$.
\tcblower
\textbf{The Trap:} The standard mathematical convention defines $\mathcal{L} = f_0 + \sum \lambda_i f_i$. Because $f_i(\mathbf{w}) \le 0$ on the feasible set and $\lambda_i \ge 0$, the product $\lambda_i f_i(\mathbf{w})$ is non-positive ($\le 0$). This guarantees that $\mathcal{L}(\mathbf{w}, \boldsymbol{\lambda}) \le f_0(\mathbf{w})$, establishing the weak duality lower bound $g(\boldsymbol{\lambda}) \le p^*$. Reversing the sign flips the direction of the inequality and invalidates the duality theorem.
\end{trapbox}

\begin{trapbox}{Gotcha 2: Assuming KKT is Sufficient for Non-Convex Problems}
An engineer solves the KKT equations for a non-convex deep learning functional and asserts that the resulting parameters must be a global minimizer.
\tcblower
\textbf{The Trap:} The KKT conditions are \textbf{necessary} for local optimality (under constraint qualifications like LICQ or Slater's condition). They are \textbf{sufficient} for global optimality if and only if the objective functional $f_0$ and all inequality constraints $f_i$ are \textbf{convex}, and all equality constraints $h_j$ are affine. On non-convex problems, points satisfying the KKT conditions can be local maxima, inflection points, or saddle points.
\end{trapbox}

\begin{trapbox}{Gotcha 3: Applying Hard-Margin SVM to Non-Linearly Separable Data}
A practitioner executes a hard-margin SVM solver on noisy empirical data containing overlapping classes.
\tcblower
\textbf{The Trap:} If the dataset is not strictly linearly separable, no pair $(\mathbf{w}, b)$ can satisfy $y_i(\mathbf{w}^T \mathbf{x}_i + b) \ge 1$ for all $i$. The primal feasible set is empty ($p^* = +\infty$). Consequently, the dual quadratic program becomes unbounded ($d^* = +\infty$), causing numerical QP solvers to crash or cycle endlessly. Non-separable data strictly requires the \textbf{Soft-Margin SVM} with slack variables $\xi_i \ge 0$ and the box constraint $0 \le \lambda_i \le C$.
\end{trapbox}

\begin{trapbox}{Gotcha 4: Forgetting to Recover the Bias Term $b^*$ in SVM}
A developer reconstructs the optimal weights $\mathbf{w}^* = \sum_{i=1}^N \lambda_i^* y_i \mathbf{x}_i$, but omits computing the intercept $b^*$.
\tcblower
\textbf{The Trap:} Without the intercept $b^*$, the decision boundary is forced to pass through the origin ($\mathbf{w}^T \mathbf{x} = 0$), destroying the maximal margin property and severely degrading accuracy. The bias $b^*$ must be recovered using any active support vector lying strictly on the margin ($0 < \lambda_k^* < C$):
\begin{equation}
b^* = y_k - \mathbf{w}^{*T} \mathbf{x}_k = y_k - \sum_{i=1}^N \lambda_i^* y_i (\mathbf{x}_i^T \mathbf{x}_k)
\end{equation}
In practical software, $b^*$ is averaged across all free support vectors to maximize numerical stability.
\end{trapbox}

\begin{trapbox}{Gotcha 5: The $N \times N$ Gram Matrix Memory Bottleneck in Kernel SVMs}
A data scientist attempts to train a Kernel SVM on a dataset with $N = 1{,}000{,}000$ samples and $d = 20$ features.
\tcblower
\textbf{The Trap:} The dual SVM requires forming or accessing the Gram kernel matrix $K_{ij} = k(\mathbf{x}_i, \mathbf{x}_j)$. Storing an $N \times N$ matrix for $N = 10^6$ requires $(10^6)^2 \times 8 \text{ bytes} = 8 \text{ Terabytes of RAM}$, causing immediate out-of-memory crashes. Dual kernel SVMs scale quadratically in storage $\mathcal{O}(N^2)$ and cubically in time $\mathcal{O}(N^3)$. When $N \gg 10^5$, linear primal SVMs (e.g., via stochastic dual coordinate ascent / LibLinear) or approximate kernel embeddings (random Fourier features) must be used.
\end{trapbox}

% ----------------------------------------------------------------------
% PART 7: COMPARATIVE SYNTHESIS & SUMMARY TABLE
% ----------------------------------------------------------------------
\section{Comparative Synthesis \& Summary Table}

Table~\ref{tab:svm_models} synthesizes the mathematical architecture of linear classification models.

\begin{table}[htbp]
\centering
\small
\begin{tabularx}{\textwidth}{lXXXX}
\toprule
\textbf{Model} & \textbf{Loss Functional} & \textbf{Regularizer} & \textbf{Sparsity in Multipliers ($\lambda_i > 0$)} & \textbf{Separability Requirement} \\
\midrule
Hard-Margin SVM & Constraint: $y_i(\mathbf{w}^T \mathbf{x}_i + b) \ge 1$ & $\frac{1}{2}\|\mathbf{w}\|_2^2$ & High (only margin boundary points) & Strictly linearly separable \\
\addlinespace
Soft-Margin SVM & Hinge Loss: $\max(0, 1 - y_i(\mathbf{w}^T\mathbf{x}_i+b))$ & $\frac{1}{2}\|\mathbf{w}\|_2^2$ & Sparse (points on or violating margin) & Handles overlapping classes ($0 \le \lambda_i \le C$) \\
\addlinespace
Logistic Regression & Log-Loss: $\log(1 + e^{-y_i(\mathbf{w}^T\mathbf{x}_i+b)})$ & $\frac{\lambda}{2}\|\mathbf{w}\|_2^2$ (optional) & Non-sparse (every sample contributes weight) & Robust to overlap; outputs probabilities \\
\addlinespace
Ridge Classification & Squared Loss: $(y_i - (\mathbf{w}^T\mathbf{x}_i+b))^2$ & $\frac{\lambda}{2}\|\mathbf{w}\|_2^2$ & Non-sparse (closed-form normal equations) & Sensitive to large positive outliers \\
\bottomrule
\end{tabularx}
\caption{Comparative Synthesis: Linear Classification and Margin Models.}
\label{tab:svm_models}
\end{table}

% ----------------------------------------------------------------------
% PART 8: PYTHON / SCIENTIFIC COMPUTING LAB
% ----------------------------------------------------------------------
\section{Python / Scientific Computing Lab}

The production script below implements the Dual Support Vector Machine Quadratic Program from scratch using \texttt{scipy.optimize.minimize}, isolates support vectors via complementary slackness, and reconstructs $\mathbf{w}^*$ and $b^*$.

\begin{lstlisting}[language=Python, caption={Dual Support Vector Machine Quadratic Programming Solver in Python.}]
import numpy as np
from scipy.optimize import minimize

def fit_dual_svm(X: np.ndarray, y: np.ndarray, C: float = 1.0):
    """Solves the Soft-Margin SVM Dual Quadratic Program:
    min 0.5 * lambda^T Q lambda - sum(lambda) subject to 0 <= lambda_i <= C and sum(lambda_i * y_i) = 0."""
    N, d = X.shape
    # Compute Gram kernel matrix Q_ij = y_i * y_j * (x_i^T x_j)
    K = X @ X.T
    Q = np.outer(y, y) * K

    def dual_objective(lam):
        return 0.5 * float(lam @ Q @ lam) - float(np.sum(lam))

    def dual_gradient(lam):
        return Q @ lam - np.ones(N)

    # Equality constraint: sum(lambda_i * y_i) = 0
    constraints = {'type': 'eq', 'fun': lambda lam: np.dot(lam, y), 'jac': lambda lam: y}
    # Box bounds: 0 <= lambda_i <= C
    bounds = [(0.0, C) for _ in range(N)]
    lam0 = np.zeros(N)

    res = minimize(dual_objective, lam0, jac=dual_gradient,
                   constraints=constraints, bounds=bounds, method='SLSQP')
    lam_opt = res.x

    # Identify support vectors (threshold above numerical precision)
    sv_mask = lam_opt > 1e-5
    sv_indices = np.where(sv_mask)[0]

    # Reconstruct weight vector w* = sum(lambda_i * y_i * x_i)
    w_star = np.sum((lam_opt * y)[:, None] * X, axis=0)

    # Recover bias b* from support vectors strictly on margin (0 < lambda_i < C)
    margin_svs = np.where((lam_opt > 1e-4) & (lam_opt < C - 1e-4))[0]
    if len(margin_svs) > 0:
        b_vals = [y[k] - float(np.dot(w_star, X[k])) for k in margin_svs]
        b_star = float(np.mean(b_vals))
    else:
        # Fallback across all active support vectors
        b_vals = [y[k] - float(np.dot(w_star, X[k])) for k in sv_indices]
        b_star = float(np.mean(b_vals)) if len(b_vals) > 0 else 0.0

    return w_star, b_star, lam_opt, sv_indices

def main():
    rng = np.random.RandomState(42)
    # Generate linearly separable synthetic 2D data
    N = 30
    X_pos = rng.randn(N // 2, 2) + np.array([2.0, 2.0])
    X_neg = rng.randn(N // 2, 2) + np.array([-2.0, -2.0])
    X = np.vstack([X_pos, X_neg])
    y = np.array([1] * (N // 2) + [-1] * (N // 2), dtype=float)

    C_val = 10.0
    w_star, b_star, lam_opt, sv_indices = fit_dual_svm(X, y, C=C_val)

    print("=================================================================")
    print(" DUAL SUPPORT VECTOR MACHINE BENCHMARK")
    print("=================================================================")
    print(f"Total Training Samples:          N = {N}")
    print(f"Box Constraint Penalty:          C = {C_val:.1f}")
    print(f"Optimal Weight Vector w*:        [{w_star[0]:.4f}, {w_star[1]:.4f}]")
    print(f"Optimal Bias Intercept b*:       {b_star:.4f}")
    print(f"Active Support Vectors Count:    {len(sv_indices)} / {N}")
    print(f"Geometric Margin Width (2/||w||): {2.0 / np.linalg.norm(w_star):.4f}")
    print("-" * 65)

    print("Active Support Vector Multipliers (lambda_i > 0):")
    for idx in sv_indices:
        print(f"  Sample {idx:<3}: label = {int(y[idx]):+d}, lambda = {lam_opt[idx]:.4f}, x = [{X[idx, 0]:.2f}, {X[idx, 1]:.2f}]")

if __name__ == "__main__":
    main()
\end{lstlisting}

% ----------------------------------------------------------------------
% PART 9: COMPREHENSIVE PRACTICE PROBLEMS
% ----------------------------------------------------------------------
\section{Comprehensive Practice Problems}

\subsection*{Part A: Analytical \& Mathematical Problems}

\begin{enumerate}[label=\textbf{P\arabic*.}]
    \item \textbf{Analytical KKT System on Quadratic Inequality:}
    Minimize $f_0(x, y) = x^2 + y^2$ subject to the inequality constraint $x + y \ge 4$ (written as $4 - x - y \le 0$).
    \begin{enumerate}[label=(\alph*)]
        \item Formulate the Lagrangian $\mathcal{L}(x, y, \lambda)$ with $\lambda \ge 0$.
        \item State all four KKT conditions.
        \item Test Case 1: $\lambda = 0$ (inactive constraint). Show that this violates primal feasibility.
        \item Test Case 2: $\lambda > 0$ (active constraint). Solve for $x^*, y^*, \lambda^*$ and verify all KKT conditions.
    \end{enumerate}

    \item \textbf{Dual of a Linear Program (LP):}
    Consider the standard primal Linear Program: $\min_{\mathbf{x} \in \mathbb{R}^n} \mathbf{c}^T \mathbf{x}$ subject to $A\mathbf{x} = \mathbf{b}$ and $\mathbf{x} \ge \mathbf{0}$.
    \begin{enumerate}[label=(\alph*)]
        \item Write the Lagrangian with dual variables $\boldsymbol{\nu} \in \mathbb{R}^m$ for equality and $\boldsymbol{\lambda} \in \mathbb{R}^n_{\ge 0}$ for non-negativity.
        \item Derive the dual function $g(\boldsymbol{\lambda}, \boldsymbol{\nu}) = \inf_{\mathbf{x}} \mathcal{L}(\mathbf{x}, \boldsymbol{\lambda}, \boldsymbol{\nu})$.
        \item State the dual Linear Program in standard maximization form.
    \end{enumerate}

    \item \textbf{Analytical 3-Point Hard-Margin SVM:}
    Consider three data points in $\mathbb{R}^2$: $\mathbf{x}_1 = [1, 1]^T$ with $y_1 = +1$, $\mathbf{x}_2 = [2, 0]^T$ with $y_2 = +1$, and $\mathbf{x}_3 = [0, 0]^T$ with $y_3 = -1$.
    \begin{enumerate}[label=(\alph*)]
        \item Set up the dual objective function $\max \sum \lambda_i - \frac{1}{2}\sum \lambda_i \lambda_j y_i y_j \mathbf{x}_i^T \mathbf{x}_j$ subject to $\lambda_1 + \lambda_2 = \lambda_3$ and $\lambda_i \ge 0$.
        \item Solve analytically for $\lambda_1^*, \lambda_2^*, \lambda_3^*$.
        \item Identify which points are support vectors and calculate the optimal separating hyperplane $(\mathbf{w}^*, b^*)$.
    \end{enumerate}

    \item \textbf{Soft-Margin SVM KKT Slack Derivation:}
    In Soft-Margin SVM with slack variables $\xi_i \ge 0$:
    \begin{equation}
    \min_{\mathbf{w}, b, \boldsymbol{\xi}} \frac{1}{2}\|\mathbf{w}\|_2^2 + C \sum_{i=1}^N \xi_i \quad \text{s.t. } y_i(\mathbf{w}^T \mathbf{x}_i + b) \ge 1 - \xi_i, \; \xi_i \ge 0
    \end{equation}
    Introduce dual multipliers $\lambda_i \ge 0$ for the margin constraint and $\mu_i \ge 0$ for the slack positivity constraint. Prove that at optimality:
    \begin{equation}
    0 \le \lambda_i \le C
    \end{equation}
    and describe the geometric location of samples when $\lambda_i = 0$, $0 < \lambda_i < C$, and $\lambda_i = C$.

    \item \textbf{Kernel Ridge Regression via Duality:}
    Consider the primal Ridge Regression problem $\min_{\mathbf{w}} \frac{1}{2}\|X\mathbf{w} - \mathbf{y}\|_2^2 + \frac{\lambda}{2}\|\mathbf{w}\|_2^2$.
    By introducing auxiliary variable $\mathbf{r} = X\mathbf{w} - \mathbf{y}$ and dual multiplier $\boldsymbol{\alpha}$, derive the dual optimization problem and show that the predictor can be written entirely in terms of the kernel matrix $K = X X^T$:
    \begin{equation}
    \hat{y}(\mathbf{x}_{\text{test}}) = \mathbf{k}(\mathbf{x}_{\text{test}})^T (K + \lambda I)^{-1} \mathbf{y}
    \end{equation}
\end{enumerate}

\subsection*{Part B: Diagnostic Multiple-Choice Questions}

\begin{enumerate}[label=\textbf{MCQ\arabic*.}]
    \item In a trained Soft-Margin Support Vector Machine with penalty parameter $C = 5.0$, a training sample has an optimal dual multiplier $\lambda_i^* = 0$. What is the geometric position of this sample relative to the decision boundary?
    \begin{enumerate}[label=(\alph*)]
        \item It lies strictly on the margin boundary $y_i(\mathbf{w}^T \mathbf{x}_i + b) = 1$.
        \item It lies strictly on the correct side of the margin: $y_i(\mathbf{w}^T \mathbf{x}_i + b) > 1$ with slack $\xi_i = 0$.
        \item It is a misclassified outlier with slack $\xi_i > 1$.
        \item It lies directly on the separating hyperplane $\mathbf{w}^T \mathbf{x}_i + b = 0$.
    \end{enumerate}

    \item Why does the Lagrange dual function $g(\boldsymbol{\lambda}, \boldsymbol{\nu}) = \inf_{\mathbf{w}} \mathcal{L}(\mathbf{w}, \boldsymbol{\lambda}, \boldsymbol{\nu})$ have the property of always being concave, even when the primal objective $f_0(\mathbf{w})$ is highly non-convex?
    \begin{enumerate}[label=(\alph*)]
        \item Because the infimum of any set of numbers is always a negative value.
        \item Because for each fixed $\mathbf{w}$, the Lagrangian is an affine (linear) function of $(\boldsymbol{\lambda}, \boldsymbol{\nu})$, and the pointwise infimum of affine functions is always concave.
        \item Because Slater's condition forces the Hessian of $g$ to be negative definite.
        \item Because the dual variables $\boldsymbol{\lambda}$ are constrained to be strictly positive.
    \end{enumerate}
\end{enumerate}
